% NI64--0C
% Insult
% 14.0
% 2013-11-30
\label{m:insult}
\EE{Taktik}: \Speech{\Ca{Vitale Bedrohung!}
Zeitkritisch (\enquote{\Speech{Time is brain.}}):
Vitalfunktionen sichern, rascher Transport an eine geeignete Abteilung.}

Grundsätzlich
ist ein Schlaganfall-Patient aufgrund seiner Diagnose als \EE{vital
bedroht} und gleichzeitig als \EE{zeitkritisch}
(\enquote{\Speech{Time is brain.}}) einzustufen. 
Eine definitive Therapie in einer Spezialabteilung
(\I{Stroke Unit}) kann andere, noch nicht abgestorbene
Teile des Hirns vor weiterem Schaden bewahren, d.\,h. es ist
sehr wichtig, dass diese Therapie möglichst schnell
erfolgt.
% 
% Bei der Behandlung des Patienten mit V.\,a. Schlaganfall
% gelten folgende Besonderheiten:


\begin{itemize}
  \item \TxMassVitMKBes
  \begin{itemize}
    \item \EE{Lagerung}: \E{leicht erhöhter Oberkörper}
    \item \EE{Sauerstoff}: Gabe je nach Sp\ce{O2},
    unterer Grenzwert 94\PC*,
    darüber ist von einer \ce{O2}-Gabe abzusehen, da es
    sonst zu einer Gefäßengstellung mit nachfolgender
    Mangeldurchblutung in den noch gesunden Hirnarealen
    kommen kann! \Commentary[Insult / Sauerstoffgabe]{In der originalen
    Lehrmeinungsnpassung ist der Grenzwert 95\PC*. Im Sinne
    der Konsistenz mit der allgemeinen Lehrmeinung ist ein
    angepasster Grenzwert von 94\PC* vertretbar
    \cite{AsboeLvWien-2013-01-29}.}
    \item \EE{Notarzt}: Grundsätzlich ist ein
    Schlaganfall-Patient aufgrund seiner Diagnose als vital
    bedroht einzustufen, er ist jedoch gleichzeitig zeitkritisch (s.\,o.).
    Daher soll von einer
    \EE{Nachberufung eines Notarztes \underline{abgesehen} werden,
    \E{wenn folgende Umstände gegeben sind}}:
    \begin{enumerate}
      \item Die Verdachtsdiagnose gilt als \E{sehr sicher}.
      \item Der Patient ist sonst als \E{stabil} (\prettyref{m:ersteinschaetzung}) eingestuft
      und
      bewusstseinsklar\footnote{Bei der Einschätzung des Bewusstseins
      darf man sich nicht durch die durch den Schlaganfall
      verursachten Symptome (Sprachstörungen, Blickstarre \ldots)
      beirren lassen.}.
      \item Es gibt sonst \E{keine anderen Umstände} die die Berufung
      eines Notarztes erforderlich machen.
      \item Es muss versucht werden, eine Aufnahme auf eine
      \E{Spezialabteilung} für Schlaganfälle zu
      organisieren\footnote{Der Erfolg ist abhängig von der
      Verfügbarkeit eines entsprechenden freien Bettes.} (Stroke Unit,
      Kontaktaufnahme mit der Leitstelle).
      \item Regionale Bestimmungen müssen beachtet werden!
    \end{enumerate}
  \end{itemize}
  \item  \EE{Blutdruck}: Blutdruckwerte bis \RRmmHg{220}{120} werden toleriert,
  wenn nicht cardiopulmonale Gründe (Myocardischämie,
  Herzinsuffizienz, \ldots) andere Grenzwerte erfordern.
  Senkung des Blutdrucks ist eine ärztliche Maßnahme.
  \item \EE{Neurocheck} inkl. Messung des \EE{Blutzuckers} ist
  zwingend erforderlich!
  \item \EE{Transportentscheidung}: Wenn verfügbar: \EE{Stroke Unit}, sonst
  Abklärung mit Leitstelle wegen Alternativen (Abt. für
  Neurologie oder Innere Medizin, mit Möglichkeit zur CT und
  Intensivüberwachung).
  Bei Anforderung einer Stroke Unit ist anzugeben:
  \begin{enumerate}
    \item Alter
    \item Beginn der Symptome
    \item Vorherige Schlaganfälle
  \end{enumerate}
\end{itemize}

\Cave{Eine engmaschige Verlaufskontrolle, insbesonders des
Bewusstseinszustandes, ist wichtig!}


\cite{Aass-1.0var1,AassMk-2013var1,MA70-Protokoll-2012-10-11,AsboeLvWien-2013-01-29,Lang2012,Ma70-Ausrueckordnung-2012-01-24}
